% Vietnamese translation for OI Public Library
% Source-File: IOI中国国家候选队论文/2005/汪汀/汪汀.ppt
\documentclass[11pt]{article}
\usepackage{oipl-vn}

\newcommand{\Oh}{\mathcal{O}}

\title{Bản trình chiếu: Ứng dụng của tìm kiếm tham số}
\author{Wang Ting}
\date{2005}

\begin{document}
\maketitle

\origtitle{Applications of parametric search}
\sourcepath{IOI China candidate team papers / 2005 / Wang Ting / Wang Ting.ppt}

\section{Phạm vi bản dịch}

Bản trình chiếu đi cùng bài viết về tìm kiếm tham số. Text stream của PPT giữ
lại các mốc \(N=9,K=3\), \(N=7,K=3\), các độ phức tạp
\(\Oh(N^2)\), \(\Oh(TN)\), \(\Oh(N\log T)\), và phép biến đổi
\[
  C_i=A_i-pB_i.
\]
Bản ghi chú này dịch mạch slide theo bài viết đã có: đưa tham số vào bài tối
ưu, biến nó thành bài quyết định, rồi dùng nhị phân hoặc lặp để tìm giá trị
tối ưu.

\section{Tư tưởng chính}

Tìm kiếm tham số không trực tiếp hỏi ``giá trị tối ưu bằng bao nhiêu''. Thay
vào đó, ta chọn một giá trị \(P\) và hỏi bài quyết định:
\[
  \text{có tồn tại nghiệm đạt ít nhất/tối đa }P\text{ hay không?}
\]
Nếu câu trả lời đơn điệu theo \(P\), ta dùng tìm kiếm nhị phân trên miền tham
số. Điểm quyết định không nằm ở phép nhị phân, mà ở hàm kiểm tra: nếu kiểm tra
giảm được từ \(\Oh(N^2)\) xuống \(\Oh(N)\), toàn bộ lời giải giảm từ quá chậm
xuống \(\Oh(N\log T)\) hoặc tương tự.

\section{Ví dụ chia đá}

Bài viết mở đầu bằng một ví dụ chia các viên đá thành các đoạn. Nếu xử lý trực
tiếp, mỗi lần thử cách chia có thể cần xét nhiều đoạn và dẫn đến độ phức tạp
lớn. Slide dùng các mốc \(N=9,K=3\) và \(N=7,K=3\) để minh họa việc đặt một
cận \(P\): ta kiểm tra có thể chia sao cho mỗi phần thỏa điều kiện liên quan
đến \(P\) hay không.

Sau khi cố định \(P\), bài toán con thường đơn giản hơn: chỉ cần quét từ trái
sang phải và đếm số đoạn hợp lệ, hoặc dùng quy hoạch động một chiều. Khi \(P\)
tăng, điều kiện khó hơn; khi \(P\) giảm, điều kiện dễ hơn. Tính đơn điệu này
cho phép nhị phân.

\section{Tối ưu tỉ số}

Mốc công thức quan trọng của PPT là biến đổi
\[
  C_i=A_i-pB_i.
\]
Nó xuất hiện khi cần tối đa hóa tỉ số
\[
  \frac{A_i+\cdots+A_j}{B_i+\cdots+B_j}.
\]
Thay vì so sánh mọi đoạn bằng tổng tiền tố trong \(\Oh(N^2)\), ta hỏi với một
tham số \(p\):
\[
  \frac{\sum A}{\sum B}\ge p
  \quad\Longleftrightarrow\quad
  \sum(A_t-pB_t)\ge 0.
\]
Như vậy bài toán tỉ số được quy về tìm đoạn có tổng \(C_t\) không âm, thường
có thể kiểm tra bằng tiền tố nhỏ nhất trong \(\Oh(N)\). Slide cũng nhắc các
hạng \(-pB_i\), \(+1-pB_i\), \(-pB_j\), cho thấy cách từng lựa chọn đóng góp
vào hàm kiểm tra sau khi đưa tham số vào.

\section{Ví dụ matching theo cận trên}

Một ví dụ khác trong bài viết là chọn một matching sao cho cạnh lớn nhất trong
matching nhỏ nhất. Với tham số \(T\), ta xóa mọi cạnh có trọng số vượt \(T\),
rồi chỉ cần hỏi đồ thị hai phía còn lại có matching hoàn hảo hay không. Nếu
có, \(T\) đủ lớn; nếu không, \(T\) quá nhỏ.

Hàm kiểm tra dùng thuật toán matching cực đại. Với mỗi \(T\), nếu kích thước
matching cực đại bằng số đỉnh cần ghép thì tồn tại phương án. Nhị phân trên
\(T\) cho giá trị nhỏ nhất. Đây là ví dụ điển hình: tham số biến một bài tối
ưu min--max thành một bài tồn tại.

\section{Kết luận}

Tìm kiếm tham số có hai mặt. Nó làm bài toán có vẻ phức tạp hơn vì phải thêm
một biến \(P\), nhưng chính biến ấy thường phá vỡ hàm mục tiêu khó xử lý và
đưa bài toán về kiểm tra tuyến tính, matching, quy hoạch động hoặc một cấu
trúc quen thuộc. Khi thiết kế lời giải, cần chứng minh tính đơn điệu của câu
trả lời và tối ưu hàm kiểm tra trước khi quan tâm đến số vòng nhị phân.

\end{document}
